\section{A King Who Gives His Life}
Christ stands before a governor who can order his death, and the Mass calls him King of the universe. How can the prisoner who declines armed rescue be the Lord to whom every ruler owes allegiance? Colossians identifies this same Son as the one through whom all things were created and through whose blood peace is made. His refusal of rescue allows the Passion to proceed; his death reconciles the people over whom he reigns.

The Introit acclaims the Lamb as slain, and the proper Preface names Christ as Priest and King in its account of his self-offering. The Mass's prayers ask for the reconciliation of divided human families and for shared heavenly reign with Christ. The Postcommunion is prayed by those already nourished by the sacrament who glory in serving him. Their present service can be understood as perseverance toward the reign they ask to share.

The first interpretation concerns how a person becomes a willing subject of this King. Augustine explains the kingdom's different origin through regeneration and identifies hearing with obedient belief; Jerome explains the Son of Man's received dominion through the Incarnation and the saints' enduring inheritance. The second concerns what the King's universal rule accomplishes among estranged creatures. Chrysostom follows the Epistle from the Son's creative primacy to reconciliation through his blood; Bellarmine explains the Psalter's universal royal claims through justice, conversion, and final judgment.

The feast celebrated here is \emph{D. N. Iesu Christi Regis}, I class, in the Sanctorale of the 1962 \emph{Missale Romanum}. Its appointed place is the last Sunday of October, which falls on 25 October in 2026. The Gloria, Creed and proper Preface belong to this occurrence. The otherwise occurring Sunday supplies no additional prayers. The eleven texts used at this Mass appear below in their liturgical order. Psalm numbers follow the Missal's Vulgate system; the usual Hebrew number appears in parentheses when different.\footnote{\emph{Missale Romanum}, Vatican typical edition, 1962, pp. 697--698, nos. 3938--3949; p. 291, no. 1080; \emph{Rubricae generales}, 17(d), 111(a), 115(c). The dated FSSP Province de France \emph{Ordo du mois}, October 2026, confirms the feast and white vestments on 25 October.}

\begin{maptable}
Introit & Apoc 5:12; 1:6; Ps 71:1 (72) & The slain Lamb receives heavenly acclaim; the psalm asks for the King's just judgment.\\
Collect & Composed oration & Sin divides human families; the prayer seeks their gathering under Christ's gentle dominion.\\
Epistle & Col 1:12--20 & Rescue and transfer lead to the Son's creative primacy and peace through the Cross.\\
Gradual & Ps 71:8,11 (72) & The King's dominion reaches all peoples and their rulers.\\
Alleluia & Dan 7:14, final clauses & The dominion of the Son of Man endures beyond destruction.\\
Gospel & John 18:33--37 & Before Pilate, Jesus identifies the kingdom's origin and the hearing that responds to truth.\\
Offertory & Ps 2:8 & The Father promises the Son the nations and the ends of the earth as his inheritance.\\
Secret & Composed oration & Sacrifice and reconciliation meet in the petition for unity and peace.\\
Preface & Proper of Christ the King & Christ is named eternal Priest and universal King in the account of his self-offering.\\
Communion & Ps 28:10--11 (29), selected clauses & The eternal King blesses his people with peace.\\
Postcommunion & Composed oration & Those already nourished and serving Christ ask to reign with him in heaven.\\
\end{maptable}

The Tract and Paschal Alleluia printed with this Mass concern seasonal votive use. They are preserved with their source notes in the terminal appendix; neither replaces a chant in the October celebration.

\tableofcontents
